% Template for the Exercisesheet - Fire Simulation I/II
% Bergische Universität Wuppertal

% !!! Compile with PDFLaTeX !!!
% Document uses the Komma-Script "article" class

\documentclass[11pt, a4paper, headheight=70pt]{scrartcl}
\voffset=2cm
%\hoffset= ... cm
%\textwidth= ... cm
%\textheight= ... cm
 
 
% Packages
%===================================
  
\usepackage[utf8]{inputenc}

\usepackage{enumitem}
	\setlist[itemize]{itemsep=4pt, topsep=0pt}
	\setlist[enumerate]{itemsep=4pt, topsep=0pt}

\usepackage{graphicx}

\usepackage{scrlayer-scrpage}

\usepackage[hidelinks]{hyperref}

\usepackage{xcolor}

\usepackage{listings}
\lstset{
xleftmargin=2em,
frame=single,
framexleftmargin=1.5em,
language=bash}

%\usepackage{showframe}


% Global Settings
%===================================

\setlength{\parindent}{0em}

\setkomafont{section}{\normalfont\bfseries\large}
\setkomafont{subsection}{\normalfont}
\setkomafont{pagehead}{\normalfont}
\setkomafont{pagefoot}{\normalfont}


% Header/Footer
%===================================

\ihead{%
	\begin{minipage}[t][63pt]{9cm}%
		\textbf{Datenanalyse mit Python}\\
		Summer Semester 2018\\[0.1cm]
		Lukas Arnold, Max Böhler
	\end{minipage}%
}

\ohead {%
	\begin{minipage}[l][65pt]{3.85cm}%
	\includegraphics[height=40pt]{BUW_Logo-schwarz.eps}%
	\end{minipage}%
}

\cfoot{\pagemark}

\ofoot{\today}


% (Re-)Newcommand
%===================================

\newcommand{\subsecline}[1]{\subsection*{#1}
\vspace{-0.5cm}
\noindent\rule{\textwidth}{0.5pt}
\vspace{0cm}}

\RedeclareSectionCommand[
  beforeskip=1cm]{subsection}



%Main Document
%===================================

\begin{document}


{\Large \textbf{Installation}}\\[0.2cm]
{\Large Python 3 und Jupyter Notebook}

\vspace{0.8cm}

Die Installation einer Python-Umgebung kann auf unterschiedlichen Wegen erfolgen. Im Rahmen dieser Veranstaltung empfehlen wir die Installation und Verwaltung von Python mittels der \textit{miniconda} Distribution. \\
Für die Überprüfung der Installation kann die Datei \texttt{installation\textunderscore test.ipynb} (siehe Moodle) mittels \textit{Jupyter Notebook} geöffnet und ausgeführt werden.


\subsecline{Linux/macOS}
\begin{enumerate}[itemsep=0.3cm]
\item Installationsdatei herunterladen:
\begin{lstlisting}[escapechar=|, basicstyle=\footnotesize]
|\href{https://conda.io/miniconda.html}{https://conda.io/miniconda.html}|
\end{lstlisting}

\item Konsole (Terminal) öffnen und im Ordner der Installationsdatei folgenden Befehl ausführen (ohne \$:):
\begin{lstlisting}[backgroundcolor = \color{lightgray}, language = Bash, framerule=0pt, basicstyle=\small]
$: bash Miniconda3-latest-Linux-x86_64.sh
\end{lstlisting}

\item Installation durchführen und Hinweistexte in der Konsole beachten.
\item Neustart der Konsole.

\item Testen der Installation durch folgenden Befehl:
\begin{lstlisting}[backgroundcolor = \color{lightgray}, language = Bash, framerule=0pt, basicstyle=\small]
$: conda --version 
\end{lstlisting}
Wird eine Miniconda-Version angezeigt, war die Installation erfolgreich.

\item Für die Installation von \textit{Jupyter Notebook} muss folgender Befehl ausgeführt \linebreak werden:
\begin{lstlisting}[backgroundcolor = \color{lightgray}, language = Bash, framerule=0pt, basicstyle=\small]
$: conda install jupyter
\end{lstlisting}

\item Installation durchführen und Hinweistexte in der Konsole beachten.
%
\item Bei erfolgreicher Installation, kann mit folgendem Befehl ein neues Notebook ge-startet werden:
\begin{lstlisting}[backgroundcolor = \color{lightgray}, language = Bash, framerule=0pt, basicstyle=\small]
$: jupyter notebook
\end{lstlisting}
\end{enumerate}

\newpage

\subsecline{Windows}

\begin{enumerate}[itemsep=0.3cm]
\item Installationsdateien herunterladen:
\begin{lstlisting}[escapechar=|, basicstyle=\footnotesize]
|\href{https://conda.io/miniconda.html}{www.conda.io/miniconda.html}|
\end{lstlisting}

\item \texttt{.exe} Datei öffnen und Installation durchführen.

\item Nach der Installation kann das Programm \textit{Anaconda Prompt} geöffnet und folgender Befehl ausgeführt werden:
\begin{lstlisting}[backgroundcolor = \color{lightgray}, language = Bash, framerule=0pt, basicstyle=\small]
$: conda --version 
\end{lstlisting}
Wird eine Miniconda-Version angezeigt, war die Installation erfolgreich.

\item Für die Installation von \textit{Jupyter Notebook} muss folgender Befehl innerhalb der \textit{Anaconda Prompt} ausgeführt werden:
\begin{lstlisting}[backgroundcolor = \color{lightgray}, language = Bash, framerule=0pt, basicstyle=\small]
$: conda install jupyter
\end{lstlisting}

\item Bei erfolgreicher Installation, kann mit folgendem Befehl innerhalb der \textit{Anaconda Prompt} ein neues Notebook geöffnet werden:
\begin{lstlisting}[backgroundcolor = \color{lightgray}, language = Bash, framerule=0pt, basicstyle=\small]
$: jupyter notebook
\end{lstlisting}

\end{enumerate}

\subsecline{Installation zusätzlicher Packete (Linux/macOS/Windows)}

Für die Installation weiterer Python Packete muss der Befehl
\begin{lstlisting}[backgroundcolor = \color{lightgray}, language = Bash, framerule=0pt, basicstyle=\small]
conda install <Packetname>
\end{lstlisting}
\vspace{0.1cm}
innerhalb des \textit{Terminals} (Linux/macOS) oder, im Fall von Windows Betriebssystemen, innerhalb der \textit{Anaconda Prompt} ausgeführt werden. Insbesondere folgende Packete werden im Rahmen dieser Veranstaltung benötigt und sollten auf dem System installiert sein:
\vspace{0.3cm}
\begin{itemize}
\setlength\itemsep{3pt}
\item \textit{matplotlib}
\item \textit{numpy}
\item \textit{pandas}
\end{itemize}


\end{document}